\section{二审程序}
\begin{enumerate}
  \item 上诉期。判决 10 日，裁定 5 日。
  \item 当事人。\bs{上诉权、辩护权}。
        \begin{enumerate}
          \item \zd{被害人不能上诉}。上诉是当事人的权利 \yjt 被告人的辩护人和近亲属需经被告人同意。
          \item 无需理由；口头、书面；找原审、上级法院都可以。
          \item 其他\zd{未上诉的同案}被告人也可以委托辩护人辩护。
        \end{enumerate}
  \item 检察院。\bs{抗诉}。
        \begin{enumerate}
          \item 抗诉。一审法院同级的检察院、确有错误、书面。\zd{民诉没有二审抗诉}。
          \item 上级认为下级\zd{抗诉不当的，先听取意见} \yjt 向同级法院撤回抗诉。
          \item 没有抗诉，不允许在二审中\bs{补充抗诉}。
          \item \zd{开庭同级检必派员}。
        \end{enumerate}
  \item 法院。\bs{撤回、2 原则、4 裁判}。
        \begin{enumerate}
          \item 撤回处理。\yi 期满前。想撤就撤。\er 期满后。\bs{死刑、判重、有错不许撤}。
          \item 审理方式。
                \begin{enumerate}
                  \item \bs{抗、死、事证、必开庭}。抗诉、死刑、事实证据有异议必开庭。
                  \item \zd{开庭后及时通知检 1 个月内阅完卷}。
                  \item 新证据应当及时通知对方。
                  \item 对同案审理中\zd{未上诉的被告人}，认为没有必要到庭的，\zd{可以不再传唤到庭}。
                  \item 要求出庭的，应当准许。
                \end{enumerate}
          \item 审判原则。
                \begin{enumerate}
                  \item \bs{全面审查原则}。就第一审判决的\zd{事实认定和法律适用}进行全面审查。
                        \al{甲盗窃杀人，一审仅判了杀人，二审不能审查盗窃}。
                  \item \bs{上诉不加刑原则}。被上被抗可加刑。
                        \begin{enumerate}
                          \item 原判罪名、罪数不当。\zd{可以改变但不得加重刑罚}。
                          \item \bs{数罪并罚只看总刑期}。
                          \item \bs{隐形加刑}。\al{盗窃 10 年 \yjt 抢劫 10 年，后者不得假释}。
                          \item 原判处罚\zd{畸轻畸重，走再审}。
                          \item 被告上诉 + 未抗诉 \yjt 被发回重审 \yjt 再抗诉，\zd{除非新的犯罪事实 + 补充起诉}，否则不得加重。
                                \al{一审 7 年 \yjt 上诉 \yjt 发回重审 \yjt 改判 3 年 \yjt 抗诉 \yjt 二审要判 \xydy 7（不得重于原审第一次）。}
                        \end{enumerate}
                \end{enumerate}
          \item 裁判。
                \begin{enumerate}
                  \item 没错\zd{维持}。
                  \item 法律错、量刑\zd{错} \yjt 应当\zd{改}判。\yi 二审法院心里有答案了。\er 提高效率。
                  \item \zd{事}实不清、\zd{证}据不足\zd{可改可发}。
                        \begin{enumerate}
                          \item \zd{事实}问题发回\zd{仅限 1 次}。
                          \item \zd{证据}问题发回\zd{不限次数}。民诉不分情况，发回仅限一次。
                          \item \zd{部分不清可以部分发回}。
                        \end{enumerate}
                  \item \zd{程序}违法 \yjt 应当\zd{发}回。次数不限。
                \end{enumerate}
          \item 审期 2 + 2(省) + X(最高法)。
          \item \zd{一经宣告即生效}。
        \end{enumerate}
  \item 发回重审。相当于要重新进行一次审理推翻先前裁判内容，因此都要撤销原判。\zd{裁定撤销原判，发回重审}。
  \item \bs{期限总结}。
        \begin{enumerate}
          \item 侦查羁押。2 + 1 + 2 + 2 + X。
          \item 审查起诉。1 + 15 日。
          \item 一审。2 + 1 + 3 + X。
          \item 二审。2 + 2 + X。
        \end{enumerate}
\end{enumerate}
